% =====================================================================
%  Método 1 — Léxico com distância de edição   (dimensão: ortografia)
%  Origem: células 3, 4, 5 e 6 do caderno
% =====================================================================
\subsection{Léxico com distância de edição}
\label{met:lexico}

O método não emprega modelo nem generalização. Opera sobre um conjunto finito de formas
atestadas, cada uma associada à sua frequência, e decompõe-se em dois passos encadeados.

O primeiro é a detecção. Dado o léxico $\lexico$,
%
\begin{equation}
\mathrm{erro}(w) = \begin{cases} 0, & w \in \lexico \\ 1, & w \notin \lexico \end{cases}
\end{equation}
%
O predicado é exato: a palavra consta ou não consta do conjunto. Constitui o único veredito
sem gradação nem limiar entre os onze métodos examinados.

O segundo é a correção. Seja $C_k(w)$ o conjunto das formas do léxico situadas a exatamente
$k$ edições de $w$. A busca é escalonada --- a distância 2 só é consultada quando a distância
1 não devolve candidato algum --- e, entre os candidatos admitidos, a frequência determina a
escolha:
%
\begin{equation}
\hat{w} = \arg\max_{c \,\in\, C(w)} f(c)
\end{equation}

A distância empregada é a de Damerau-Levenshtein, que computa a transposição de dois
caracteres adjacentes como uma única edição. A escolha não é detalhe de implementação: é ela
que situa \texttt{nativa} a uma edição de \texttt{natvia}. Sob métrica que não contemple a
transposição, a permuta custaria 2, \texttt{nativa} ficaria fora do conjunto de candidatos e
a correção emitida seria \texttt{naevia}.

O algoritmo não busca formas semelhantes. Ele enumera as 949 sequências situadas a uma
edição de \texttt{natvia} e submete cada uma ao teste de pertinência a $\lexico$; quatro
sobrevivem, conforme a Tabela~\ref{tab:candidatos}.

\begin{table}[H]
\centering
\dimtable{
\begin{tabular}{@{}lll r@{}}
\toprule
$c$ & Operação sobre \texttt{natvia} & & $f(c)$ \\
\midrule
\texttt{nativa} & transposição de \texttt{v} e \texttt{i} & & 912 \\
\texttt{naevia} & substituição de \texttt{t} por \texttt{e}  & & 103 \\
\texttt{natia}  & remoção do \texttt{v}              & &  77 \\
\texttt{latvia} & substituição de \texttt{n} por \texttt{l}  & &  23 \\
\bottomrule
\end{tabular}}
\caption{Candidatos a uma edição de \texttt{natvia}, com sua frequência no léxico.}
\label{tab:candidatos}
\end{table}

Os quatro candidatos equidistam da forma detectada, de modo que a decisão recai integralmente
sobre $f$. A frequência não opera como critério auxiliar de desempate, mas como único
critério.
